\begin{pastexam}[label={px:nccu-113-dsa-01}]{政大資科所 113 學年度 第 1 題}
（單選；不必說明理由）

(4\%) Which of the following statements is wrong:
\begin{enumerate}[label=\Alph*.]
\item $n\log n=\Omega(n^{1.001})$
\item $\frac{1}{1}+\frac{1}{2}+\frac{1}{3}+\frac{1}{4}+\frac{1}{5}+\cdots+\frac{1}{n}=O(\log n)$
\item $\log(n!)=\Theta(n\log n)$
\item $1+2+4+8+16+\cdots+2^k=\Theta(2^k)$
\end{enumerate}
\end{pastexam}

\begin{pastsoln}
答案是 A。 題目問哪一個是\textbf{錯的}。

\begin{itemize}
\item[A.] 錯。 任何正的次方都比對數長得快（資結書第~1.6~節）：
\[
\frac{n \log n}{n^{1.001}} = \frac{\log n}{n^{0.001}} \to 0 \quad (n \to \infty),
\]
所以 $n \log n = o(n^{1.001})$，不可能是 $\Omega(n^{1.001})$。 要注意 $n^{0.001}$ 追上 $\log n$ 的 $n$ 非常大
（以 $2$ 為底，大約要 $n > 2^{13747}$），小的 $n$ 看起來 $n \log n$ 比較大，但漸近符號只看 $n$ 夠大的時候。
\item[B.] 對。 這是調和數 $H_n$。 用積分夾（資結書定理~A.7）：$H_n \le 1 + \int_1^n \frac{dx}{x} = 1 + \ln n$，
所以 $H_n = O(\log n)$（其實 $H_n = \Theta(\log n)$）。
\item[C.] 對。 上界：$n! \le n^n$，$\log(n!) \le n \log n$。 下界：$n!$ 的後一半因數 $\lceil n/2\rceil, \ldots, n$
每個都 $\ge n/2$，所以 $n! \ge (n/2)^{n/2}$，$\log(n!) \ge \frac{n}{2}\log\frac{n}{2} = \Omega(n \log n)$。
\item[D.] 對。 等比級數 $1 + 2 + \cdots + 2^k = 2^{k+1} - 1$（資結書定理~A.4），夾在 $2^k$ 和 $2 \cdot 2^k$ 之間，
是 $\Theta(2^k)$。
\end{itemize}
\end{pastsoln}
